problem:  
Let \((M^n, g, f)\) be a complete, noncompact **shrinking gradient Ricci soliton** satisfying  
\[
\mathrm{Ric}_g + \nabla^2 f = \frac{1}{2} g,
\]
normalized so that \(f(p)=0\) at some base point \(p\in M\).  
Assume that \(M\) has **Euclidean volume growth** and **bounded nonnegative curvature operator**.  
Known examples suggest that the asymptotic geometry of such a soliton can exhibit either cylindrical or conical behavior — e.g., the Bryant soliton is asymptotically paraboloidal rather than conical.  

**Problem:**  
In dimensions \(n \le 5\), does there exist a **nontrivial asymptotic gap phenomenon** for such solitons?  
Namely, is it true that any sequence of rescalings  
\[
(M, r_i^{-2}g, p)
\]
with \(r_i \to \infty\) converges (in the Cheeger–Gromov sense) to a **unique tangent cone at infinity**, and moreover, that the link of this tangent cone is a smooth Einstein manifold?  
In particular, is the asymptotic cone for 3‑ or 4‑dimensional shrinkers always smooth (not singular) under these curvature and volume assumptions?

Why is it a "good" problem:  
This formulation targets a central, **currently unresolved** question in shrinking Ricci soliton theory: the **asymptotic rigidity and uniqueness of tangent cones** for noncompact shrinkers.  
Unlike a direct classification statement, it addresses the **analytic and geometric mechanism** governing the large‑scale limits of solitons. The question remains open even in low dimensions and under strong curvature conditions.  
A positive answer would connect PDE analysis (long‑distance behavior of the soliton potential), Riemannian geometry (structure of noncollapsed limits with bounded curvature), and geometric measure theory (regularity of metric cones).  
It is simple to state yet deeply interwoven with existing soliton classification and Ricci flow stability results, and provides a natural bridge between local Ricci‑flow regularity theory and global geometric convergence phenomena.